\chapter{Membership}

\section{Types of Membership}

\subsection{Memberships Offers}
The Association shall offer three (3) types of memberships that are not 
necessarily mutually exclusive:
\begin{enumerate}

    \item General (voting) membership,

    \item Honorary (non-voting) membership, and

    \item Paid (non-voting) membership.

\end{enumerate}

\subsection{Terms of Membership}
The term of all memberships shall start after May 1st and end on April 
30th of the following year.

\subsection{Membership List}
The Association shall maintain a list of members of the categories
and ensure that membership is being offered in accordance with
this constitution for every academic year.

\section{General Membership}

\subsection{General Membership Eligibility}
The eligible for general (voting) membership shall be all undergraduate
students currently registered at the CMS department.

\subsection{General Membership Termination}
General (voting) membership of the Association shall be 
automatically terminated if a member ceases to be a registered
student at the CMS department or withdraws their membership.

\subsection{General Membership Fee}
There shall be no fee for general (voting) membership.


\subsection{General Membership Rights}
General members shall have the rights to:
\begin{enumerate}

    \item Have one vote in Association referenda, and
    Members Meetings, and

    \item Sign petitions related to the Association to
    enact referenda or to call for a special meeting 
    of the members.

\end{enumerate}

\section{Honorary Membership}

\subsection{Honorary Membership Purpose}
The purpose of honorary membership is to recognize
and honour individuals who have demonstrated exceptional dedication
and service to the Association or the CMS department. The membership
is intended to acknowledge faculty, staff, alumni, or community members
who have positively impacted the Association or the CMS students.

\subsection{Honorary Membership Eligibility}
Honorary membership may be granted to any individual 
who has made significant contributions to the Association or the 
CMS department. They will be nominated and approved by the Board of 
Executives.

\subsection{Honorary Membership Termination}
Honorary membership may be revoked by a two-thirds
(2/3) majority vote of the Board of Executives. Honorary members
may also voluntarily withdraw their membership at any time.

\section{Paid Membership}

\subsection{Paid Membership Purpose}
The purpose of paid membership is to provide additional
benefits and services to members who choose to support the Association
financially. Paid members may receive exclusive access to certain events,
resources, or discounts, as determined by the Board of
Executives.

\subsection{Paid Membership Eligibility}
Paid membership may be offered to individuals who may
already be general members or honorary members,
or to individuals who do not meet the criteria for general or honorary
membership.

\subsection{Paid Membership Fee}
The fee for paid membership shall be determined by the
Board of Executives and may be subject to change on an annual basis.

\subsection{Paid Membership Termination}
Paid membership shall be automatically terminated if a
member fails to pay the membership fee within the specified time frame
or withdraws their membership.

\subsection{Paid Membership Rights}
Paid members shall have the rights to access the benefits
and services associated with their membership, as determined by the
Board of Executives for the academic year.
However, they shall not have voting rights unless they are also
a general member.
